% Silabus Mata Kuliah Kriptografi - OBE
% Program Studi Matematika S1
% Format Akademik Lengkap

\documentclass[12pt,a4paper]{article}
\usepackage[utf8]{inputenc}
\usepackage[T1]{fontenc}
\usepackage[indonesian]{babel}
\usepackage[verbose=false]{geometry}
\usepackage{booktabs}
\usepackage{longtable}
\usepackage{array}
\usepackage{enumitem}
\usepackage{listings}
\usepackage{xcolor}
\usepackage{hyperref}
\usepackage{graphicx}
\usepackage{fancyhdr}

\geometry{margin=2.5cm}
\setlength{\parindent}{0pt}
\setlength{\parskip}{6pt}

% Listings style
\lstset{
    basicstyle=\ttfamily\small,
    breaklines=true,
    frame=single,
    numbers=left,
    numberstyle=\tiny,
    keywordstyle=\color{blue},
    commentstyle=\color{green!60!black},
    stringstyle=\color{red},
    showstringspaces=false,
    tabsize=2
}

\lstdefinestyle{matlab}{
    language=Matlab,
    morekeywords={function,end,if,else,for,while,return}
}

\lstdefinestyle{cstyle}{
    language=C,
    morekeywords={int,char,void,return,for,while,if,else}
}

\lstdefinestyle{cppstyle}{
    language=C++,
    morekeywords={int,char,void,return,for,while,if,else,string,cout,cin}
}

\lstdefinestyle{pascalstyle}{
    language=Pascal,
    morekeywords={program,var,begin,end,function,procedure,if,then,else,for,to,do,while}
}

\title{\textbf{SILABUS MATA KULIAH KRIPTOGRAFI}\\
\large Berbasis Outcome-Based Education (OBE)\\
Program Studi Matematika -- S1}
\author{Departemen Matematika}
\date{}

\begin{document}

\maketitle
\thispagestyle{empty}
\newpage

\tableofcontents
\newpage

%==============================================================================
\section{Identitas Mata Kuliah}
%==============================================================================

\begin{table}[!htbp]
\centering
\small
\begin{tabular}{|l|l|}
\hline
\textbf{Aspek} & \textbf{Keterangan} \\
\hline
Nama Mata Kuliah & Kriptografi \\
Fakultas & Fakultas MIPA / Sains \\
Program Studi & Matematika \\
Kode & MAT-XXX \\
Bobot SKS & 3 SKS \\
Kelompok & Program Studi \\
Sifat Pengambilan & Wajib \\
Semester Ke & Ganjil/Genap (sesuai kurikulum) \\
Ketersediaan & Terbuka \\
Bentuk Pembelajaran & Kelas, Praktikum Laboratorium \\
Media & Campuran (Blended Learning) \\
Rumpun Mata Kuliah & KBK Matematika Terapan / Computational \\
Prasyarat & Aljabar Linear, Matematika Diskrit \\
\hline
\end{tabular}
\caption{Identitas Mata Kuliah Kriptografi (format BAN-PT/LAM)}
\end{table}

\textit{Catatan:} Jumlah pertemuan 16 minggu; komposisi pembelajaran 30\% teori, 70\% praktik.

%==============================================================================
\section{Deskripsi Mata Kuliah}
%==============================================================================

Mata kuliah Kriptografi dalam Kurikulum Program Studi Matematika diberikan kepada mahasiswa semester yang ditetapkan dengan bobot 3 SKS. Mata kuliah ini merupakan mata kuliah wajib yang ditempuh dengan prasyarat Aljabar Linear dan Matematika Diskrit. Mata kuliah ini bertujuan untuk mendukung Capaian Pembelajaran Lulusan (CPL) berupa kemampuan menguasai konsep teoritis matematika dan komputasi, menerapkan pemikiran logis dan inovatif, serta merancang dan mengimplementasikan algoritma matematika.

Materi mencakup prinsip-prinsip matematika kriptografi (aritmetika modular, algoritma Euclid, teori bilangan), algoritma kriptografi klasik dan modern (Caesar, Vigenère, DES, AES, RSA), fungsi hash, tanda tangan digital, protokol kriptografi, manajemen kunci, dan kriptanalisis dasar. Pembelajaran menekankan pemahaman matematis dan implementasi algoritma menggunakan MATLAB, C, C++, dan Pascal, dengan pendekatan Outcome-Based Education (OBE), project-based learning, dan praktikum laboratorium.

%==============================================================================
\section{Capaian Pembelajaran Lulusan (CPL) yang Relevan}
%==============================================================================

\begin{table}[!htbp]
\centering
\small
\begin{tabular}{|c|>{\raggedright\arraybackslash}p{0.72\linewidth}|}
\hline
\textbf{Kode CPL} & \textbf{Rumusan CPL} \\
\hline
CPL-1 & Pengetahuan: Menguasai konsep teoritis matematika termasuk aljabar, analisis, dan matematika diskrit untuk menyelesaikan masalah. \\
\hline
CPL-2 & Pengetahuan: Menguasai prinsip-prinsip komputasi dan pemrograman untuk implementasi solusi matematika. \\
\hline
CPL-3 & Keterampilan Umum: Mampu menerapkan pemikiran logis, kritis, sistematis, dan inovatif dalam konteks pengembangan ilmu pengetahuan. \\
\hline
CPL-4 & Keterampilan Khusus: Mampu merancang dan mengimplementasikan algoritma matematika menggunakan bahasa pemrograman. \\
\hline
CPL-5 & Keterampilan Khusus: Mampu menganalisis dan menyelesaikan masalah matematika terapan. \\
\hline
\end{tabular}
\caption{CPL Program Studi Matematika yang dibebankan pada mata kuliah Kriptografi}
\end{table}

%==============================================================================
\section{Capaian Pembelajaran Mata Kuliah (CPMK)}
%==============================================================================

\begin{table}[!htbp]
\centering
\small
\begin{tabular}{|>{\raggedright\arraybackslash}p{3.4cm}|c|>{\raggedright\arraybackslash}p{6.2cm}|}
\hline
\textbf{Kode CPL yang Didukung} & \textbf{Kode CPMK} & \textbf{Rumusan CPMK} \\
\hline
CPL-1, CPL-3 & CPMK-1 & Menjelaskan prinsip-prinsip dasar kriptografi dan konsep keamanan informasi. \\
\hline
CPL-1, CPL-3, CPL-4, CPL-5 & CPMK-2 & Menerapkan aritmetika modular, algoritma Euclid, dan teori bilangan dalam konteks kriptografi. \\
\hline
CPL-1, CPL-2, CPL-3, CPL-4, CPL-5 & CPMK-3 & Menganalisis dan mengimplementasikan algoritma kriptografi klasik (Caesar, Vigenère, substitusi). \\
\hline
CPL-1, CPL-2, CPL-3, CPL-4, CPL-5 & CPMK-4 & Menganalisis dan mengimplementasikan algoritma kriptografi modern simetris (AES, DES) dan asimetris (RSA). \\
\hline
CPL-1, CPL-2, CPL-3, CPL-4, CPL-5 & CPMK-5 & Menerapkan fungsi hash dan tanda tangan digital dalam keamanan data. \\
\hline
CPL-1, CPL-3, CPL-4, CPL-5 & CPMK-6 & Merancang dan menganalisis protokol kriptografi serta manajemen kunci. \\
\hline
CPL-1, CPL-3, CPL-4, CPL-5 & CPMK-7 & Melakukan kriptanalisis dasar terhadap cipher sederhana. \\
\hline
CPL-2, CPL-3, CPL-4, CPL-5 & CPMK-8 & Mengimplementasikan algoritma kriptografi menggunakan MATLAB, C, C++, dan Pascal. \\
\hline
\end{tabular}
\caption{Capaian Pembelajaran Mata Kuliah (CPMK) Kriptografi}
\end{table}

%==============================================================================
\section{Sub-CPMK dan Indikator Kinerja}
%==============================================================================

\subsection{Sub-CPMK 1: Prinsip Dasar Kriptografi}
\begin{itemize}
\item \textbf{Indikator 1.1:} Menjelaskan model komunikasi aman dan ancaman keamanan.
\item \textbf{Indikator 1.2:} Mendefinisikan confidentiality, integrity, dan authenticity.
\item \textbf{Indikator 1.3:} Membedakan kriptografi simetris dan asimetris.
\end{itemize}

\subsection{Sub-CPMK 2: Matematika Kriptografi}
\begin{itemize}
\item \textbf{Indikator 2.1:} Menghitung operasi aritmetika modular (penjumlahan, perkalian, invers).
\item \textbf{Indikator 2.2:} Mengimplementasikan algoritma Euclid dan Euclid diperluas.
\item \textbf{Indikator 2.3:} Menentukan bilangan prima dan uji primalitas.
\item \textbf{Indikator 2.4:} Menghitung pangkat modular (fast exponentiation).
\end{itemize}

\subsection{Sub-CPMK 3: Algoritma Kriptografi Klasik}
\begin{itemize}
\item \textbf{Indikator 3.1:} Mengimplementasikan cipher Caesar dan Vigenère.
\item \textbf{Indikator 3.2:} Menganalisis kelemahan cipher klasik.
\item \textbf{Indikator 3.3:} Melakukan enkripsi dan dekripsi dengan cipher substitusi.
\end{itemize}

\subsection{Sub-CPMK 4: Algoritma Kriptografi Modern}
\begin{itemize}
\item \textbf{Indikator 4.1:} Menjelaskan struktur AES dan DES.
\item \textbf{Indikator 4.2:} Mengimplementasikan algoritma RSA.
\item \textbf{Indikator 4.3:} Membedakan mode operasi block cipher (ECB, CBC).
\end{itemize}

\subsection{Sub-CPMK 5: Fungsi Hash dan Tanda Tangan Digital}
\begin{itemize}
\item \textbf{Indikator 5.1:} Menjelaskan sifat fungsi hash kriptografis.
\item \textbf{Indikator 5.2:} Mengimplementasikan SHA/MD5 untuk integritas data.
\item \textbf{Indikator 5.3:} Merancang tanda tangan digital berbasis RSA.
\end{itemize}

\subsection{Sub-CPMK 6: Protokol dan Manajemen Kunci}
\begin{itemize}
\item \textbf{Indikator 6.1:} Menganalisis protokol Diffie-Hellman.
\item \textbf{Indikator 6.2:} Menjelaskan siklus hidup manajemen kunci.
\item \textbf{Indikator 6.3:} Merancang protokol pertukaran kunci yang aman.
\end{itemize}

\subsection{Sub-CPMK 7: Kriptanalisis}
\begin{itemize}
\item \textbf{Indikator 7.1:} Melakukan frequency analysis pada cipher substitusi.
\item \textbf{Indikator 7.2:} Menyerang cipher Vigenère dengan Kasiski examination.
\item \textbf{Indikator 7.3:} Mengevaluasi kekuatan algoritma terhadap serangan.
\end{itemize}

\subsection{Sub-CPMK 8: Implementasi Pemrograman}
\begin{itemize}
\item \textbf{Indikator 8.1:} Mengimplementasikan cipher dasar dalam MATLAB.
\item \textbf{Indikator 8.2:} Mengimplementasikan cipher dalam C dan C++.
\item \textbf{Indikator 8.3:} Mengimplementasikan cipher dalam Pascal.
\item \textbf{Indikator 8.4:} Membuat dokumentasi dan pengujian kode.
\end{itemize}

%==============================================================================
\section{Peta Keterkaitan CPMK terhadap CPL}
%==============================================================================

\begin{table}[!htbp]
\centering
\small
\begin{tabular}{|c|c|c|c|c|c|}
\hline
\textbf{CPMK} & \textbf{CPL-1} & \textbf{CPL-2} & \textbf{CPL-3} & \textbf{CPL-4} & \textbf{CPL-5} \\
\hline
CPMK-1 & $\bullet$ & & $\bullet$ & & \\
CPMK-2 & $\bullet$ & & $\bullet$ & $\bullet$ & $\bullet$ \\
CPMK-3 & $\bullet$ & $\bullet$ & $\bullet$ & $\bullet$ & $\bullet$ \\
CPMK-4 & $\bullet$ & $\bullet$ & $\bullet$ & $\bullet$ & $\bullet$ \\
CPMK-5 & $\bullet$ & $\bullet$ & $\bullet$ & $\bullet$ & $\bullet$ \\
CPMK-6 & $\bullet$ & & $\bullet$ & $\bullet$ & $\bullet$ \\
CPMK-7 & $\bullet$ & & $\bullet$ & $\bullet$ & $\bullet$ \\
CPMK-8 & & $\bullet$ & $\bullet$ & $\bullet$ & $\bullet$ \\
\hline
\end{tabular}
\caption{Pemetaan CPMK terhadap CPL (simbol $\bullet$ = berkontribusi)}
\end{table}

%==============================================================================
\section{Bahan Kajian dan Referensi Utama}
%==============================================================================

\begin{table}[!htbp]
\centering
\small
\begin{tabular}{|p{5.5cm}|p{7.5cm}|}
\hline
\textbf{Bahan Kajian} & \textbf{Referensi Utama} \\
\hline
Matematika kriptografi: aritmetika modular, algoritma Euclid, teori bilangan, fast exponentiation & Stinson \& Paterson (2018). \textit{Cryptography: Theory and Practice} (4th ed.). CRC Press. \\
\hline
Algoritma kriptografi klasik dan modern: Caesar, Vigenère, DES, AES, RSA; fungsi hash, tanda tangan digital & Katz \& Lindell (2020). \textit{Introduction to Modern Cryptography} (3rd ed.). CRC Press. \\
\hline
Protokol kriptografi, manajemen kunci, kriptanalisis; implementasi pemrograman (MATLAB, C, C++, Pascal) & Trappe \& Washington (2006). \textit{Introduction to Cryptography with Coding Theory} (2nd ed.). Pearson. \\
\hline
\end{tabular}
\caption{Bahan Kajian dan Referensi Utama (format BAN-PT/LAM)}
\end{table}

%==============================================================================
\section{Rencana Pembelajaran Semester (RPS)}
%==============================================================================

RPS dirancang untuk 16 minggu dengan estimasi waktu per minggu: PB (proses belajar tatap muka), PT (penugasan terstruktur), BM (belajar mandiri). Total beban 3 SKS $\approx$ 150 menit/minggu. Bentuk pembelajaran: kelas dan praktikum laboratorium; media campuran (blended learning).

\begin{longtable}{|c|c|c|>{\raggedright\arraybackslash}p{2.4cm}|>{\raggedright\arraybackslash}p{2.5cm}|>{\raggedright\arraybackslash}p{2.1cm}|}
\hline
\textbf{Mg} & \textbf{Tanggal} & \textbf{Sub-CPMK} & \textbf{Materi / Bahan kajian} & \textbf{Metode pembelajaran} & \textbf{PB / PT / BM} \\
\hline
\endfirsthead

\hline
\textbf{Mg} & \textbf{Tanggal} & \textbf{Sub-CPMK} & \textbf{Materi / Bahan kajian} & \textbf{Metode pembelajaran} & \textbf{PB / PT / BM} \\
\hline
\endhead

1 & Rencana semester & Sub-CPMK 1 & Pengantar Kriptografi, model keamanan, ancaman & Kelas + diskusi kasus & 45 / 30 / 75 \\
\hline
2 & Rencana semester & Sub-CPMK 2 & Aritmetika modular, operasi dasar, invers modular & Kelas + praktikum MATLAB & 45 / 30 / 75 \\
\hline
3 & Rencana semester & Sub-CPMK 2 & Algoritma Euclid dan Euclid diperluas & Kelas + praktikum C & 45 / 30 / 75 \\
\hline
4 & Rencana semester & Sub-CPMK 2 & Teori bilangan: prima, fast exponentiation & Kelas + praktikum C++ & 45 / 30 / 75 \\
\hline
5 & Rencana semester & Sub-CPMK 3 & Cipher Caesar dan substitusi sederhana & Kelas + lab (4 bahasa) & 30 / 45 / 75 \\
\hline
6 & Rencana semester & Sub-CPMK 3 & Cipher Vigenère dan polyalphabetic & Kelas + lab Vigenère & 30 / 45 / 75 \\
\hline
7 & Rencana semester & Sub-CPMK 4 & Block cipher: DES, AES, mode ECB/CBC & Kelas + simulasi/demo & 45 / 30 / 75 \\
\hline
8 & Rencana semester & -- & \textbf{Ujian Tengah Semester} & UTS tertulis & 100 / 0 / 50 \\
\hline
9 & Rencana semester & Sub-CPMK 4 & Algoritma RSA, enkripsi asimetris & Kelas + implementasi RSA & 45 / 30 / 75 \\
\hline
10 & Rencana semester & Sub-CPMK 5 & Fungsi hash (SHA, MD5), integritas data & Kelas + lab hash & 30 / 45 / 75 \\
\hline
11 & Rencana semester & Sub-CPMK 5 & Tanda tangan digital, implementasi RSA-based & Kelas + praktikum & 30 / 45 / 75 \\
\hline
12 & Rencana semester & Sub-CPMK 6 & Protokol kriptografi, Diffie-Hellman, TLS & Kelas + analisis protokol & 45 / 30 / 75 \\
\hline
13 & Rencana semester & Sub-CPMK 6 & Manajemen kunci, siklus hidup, distribusi & Kelas + studi kasus & 30 / 45 / 75 \\
\hline
14 & Rencana semester & Sub-CPMK 7 & Kriptanalisis, frequency analysis, Kasiski & Kelas + lab kriptanalisis & 30 / 45 / 75 \\
\hline
15 & Rencana semester & Sub-CPMK 8 & Presentasi proyek, implementasi multi-bahasa & Kelas + demo proyek & 30 / 45 / 75 \\
\hline
16 & Rencana semester & -- & \textbf{Ujian Akhir Semester} & UAS tertulis & 100 / 0 / 50 \\
\hline

\end{longtable}

\textit{Catatan: Mg = Minggu ke. PB = proses belajar (menit), PT = penugasan terstruktur (menit), BM = belajar mandiri (menit). Tanggal diisi sesuai jadwal semester.}

%==============================================================================
\section{Strategi Pembelajaran OBE}
%==============================================================================

\subsection{Project-Based Learning (PBL)}
Mahasiswa menyelesaikan proyek implementasi algoritma kriptografi dalam minimal dua bahasa pemrograman (MATLAB, C, C++, atau Pascal). Proyek mencakup dokumentasi, pengujian, dan presentasi.

\subsection{Problem-Based Learning}
Studi kasus keamanan informasi dianalisis untuk mengidentifikasi kebutuhan kriptografi dan merancang solusi. Diskusi kelompok mendorong pemikiran kritis.

\subsection{Praktikum/Laboratorium}
Setiap pertemuan praktik mencakup:
\begin{itemize}
\item Implementasi algoritma sesuai topik.
\item Pengujian dengan data sampel.
\item Analisis hasil dan pembahasan kelemahan.
\item Laporan singkat praktikum.
\end{itemize}

%==============================================================================
\section{Metode Penilaian dan Rubrik}
%==============================================================================

Penilaian hasil belajar mengacu pada Capaian Pembelajaran Mata Kuliah (CPMK) dan Sub-CPMK yang tercantum dalam matriks RPS di atas.

\subsection{Komponen Penilaian}
\begin{table}[!htbp]
\centering
\begin{tabular}{|l|l|l|}
\hline
\textbf{Komponen} & \textbf{Bobot} & \textbf{Bentuk} \\
\hline
Ujian Tengah Semester (UTS) & 25\% & Tertulis \\
Ujian Akhir Semester (UAS) & 35\% & Tertulis \\
Tugas/Praktikum & 25\% & Laporan + Kode \\
Proyek Implementasi & 15\% & Demo + Dokumentasi \\
\hline
\end{tabular}
\caption{Komponen Penilaian}
\end{table}

\subsection{Rubrik Penilaian Praktikum (Skala 1--4)}
\begin{table}[!htbp]
\centering
\small
\begin{tabular}{|>{\raggedright\arraybackslash}p{2.2cm}|>{\raggedright\arraybackslash}p{2.5cm}|>{\raggedright\arraybackslash}p{2.5cm}|>{\raggedright\arraybackslash}p{2.5cm}|}
\hline
\textbf{Kriteria} & \textbf{4 (Sangat Baik)} & \textbf{3 (Baik)} & \textbf{2 (Cukup)} \\
\hline
Kebenaran Algoritma & Kode benar, output sesuai & Kode benar, minor bug & Kode sebagian benar \\
\hline
Dokumentasi & Lengkap, jelas & Cukup lengkap & Kurang lengkap \\
\hline
Pengujian & Test cases lengkap & Beberapa test cases & Minimal testing \\
\hline
Analisis & Analisis mendalam & Analisis memadai & Analisis sederhana \\
\hline
\end{tabular}
\caption{Rubrik Penilaian Praktikum}
\end{table}

\subsection{Rubrik Penilaian Proyek}
\begin{table}[!htbp]
\centering
\small
\begin{tabular}{|>{\raggedright\arraybackslash}p{2.5cm}|>{\raggedright\arraybackslash}p{1.9cm}|>{\raggedright\arraybackslash}p{1.9cm}|>{\raggedright\arraybackslash}p{1.9cm}|>{\raggedright\arraybackslash}p{1.9cm}|}
\hline
\textbf{Aspek} & \textbf{4} & \textbf{3} & \textbf{2} & \textbf{1} \\
\hline
Implementasi & $\geq$3 bahasa, lengkap & 2 bahasa, lengkap & 2 bahasa, partial & 1 bahasa \\
\hline
Fungsionalitas & Semua fitur jalan & Mayoritas jalan & Sebagian jalan & Tidak jalan \\
\hline
Presentasi & Jelas, interaktif & Jelas & Cukup jelas & Tidak jelas \\
\hline
Dokumentasi & Profesional & Memadai & Minimal & Tidak ada \\
\hline
\end{tabular}
\caption{Rubrik Penilaian Proyek}
\end{table}

%==============================================================================
\section{Daftar Topik Kriptografi}
%==============================================================================

\begin{enumerate}
\item \textbf{Matematika Kriptografi}
\begin{itemize}
\item Aritmetika modular
\item Algoritma Euclid dan Euclid diperluas
\item Teori bilangan (prima, Fermat, Euler)
\item Pangkat modular (fast exponentiation)
\end{itemize}

\item \textbf{Algoritma Kriptografi Klasik}
\begin{itemize}
\item Cipher Caesar, substitusi monoalfabetik
\item Cipher Vigenère, polyalphabetic
\item Transposisi
\end{itemize}

\item \textbf{Algoritma Kriptografi Modern}
\begin{itemize}
\item Kriptografi simetris: DES, AES
\item Kriptografi asimetris: RSA
\item Mode operasi block cipher
\end{itemize}

\item \textbf{Protokol Kriptografi}
\begin{itemize}
\item Protokol pertukaran kunci (Diffie-Hellman)
\item Protokol autentikasi
\item TLS/SSL
\end{itemize}

\item \textbf{Kriptanalisis}
\begin{itemize}
\item Frequency analysis
\item Kasiski examination
\item Serangan known-plaintext, chosen-plaintext
\end{itemize}

\item \textbf{Aplikasi Kriptografi}
\begin{itemize}
\item Fungsi hash (SHA, MD5)
\item Tanda tangan digital
\item Manajemen kunci
\item Keamanan komunikasi
\end{itemize}
\end{enumerate}

%==============================================================================
\section{Contoh Tugas Proyek dan Studi Kasus}
%==============================================================================

\subsection{Tugas Proyek 1: Implementasi Multi-Bahasa}
Implementasikan cipher Caesar dalam empat bahasa pemrograman (MATLAB, C, C++, Pascal). Bandingkan performa dan kemudahan implementasi. Buat dokumentasi penggunaan.

\subsection{Tugas Proyek 2: Sistem Enkripsi Hibrida}
Rancang sistem yang menggabungkan AES (simetris) dan RSA (asimetris) untuk enkripsi file. Implementasikan dalam minimal dua bahasa. Jelaskan alur enkripsi dan dekripsi.

\subsection{Studi Kasus 1: Keamanan Pesan Instan}
Analisis bagaimana aplikasi pesan instan (WhatsApp, Signal) menggunakan kriptografi end-to-end. Identifikasi algoritma yang digunakan dan manajemen kunci.

\subsection{Studi Kasus 2: Tanda Tangan Dokumen}
Sebuah universitas ingin menerapkan tanda tangan digital untuk surat resmi. Rancang solusi berbasis RSA dan hash. Jelaskan protokol verifikasi.

%==============================================================================
\section{Contoh Kode Kriptografi}
%==============================================================================

\subsection{MATLAB: Cipher Caesar}
\begin{lstlisting}[style=matlab,caption={Implementasi Caesar Cipher dalam MATLAB}]
function ciphertext = caesar_encrypt(plaintext, shift)
% Caesar Cipher Encryption
% plaintext: string input
% shift: integer (0-25)
    plaintext = upper(plaintext);
    ciphertext = '';
    for i = 1:length(plaintext)
        if plaintext(i) >= 'A' && plaintext(i) <= 'Z'
            c = double(plaintext(i)) - 65;
            c = mod(c + shift, 26);
            ciphertext = [ciphertext char(c + 65)];
        else
            ciphertext = [ciphertext plaintext(i)];
        end
    end
end

function plaintext = caesar_decrypt(ciphertext, shift)
    plaintext = caesar_encrypt(ciphertext, 26 - shift);
end
\end{lstlisting}

\subsection{C: Cipher Caesar}
\begin{lstlisting}[style=cstyle,caption={Implementasi Caesar Cipher dalam C}]
#include <stdio.h>
#include <ctype.h>
#include <string.h>

void caesar_encrypt(char *plaintext, int shift) {
    int i;
    for (i = 0; plaintext[i] != '\0'; i++) {
        if (isalpha(plaintext[i])) {
            char base = isupper(plaintext[i]) ? 'A' : 'a';
            plaintext[i] = base + (plaintext[i] - base + shift) % 26;
        }
    }
}

void caesar_decrypt(char *ciphertext, int shift) {
    caesar_encrypt(ciphertext, 26 - shift);
}

int main() {
    char msg[] = "HELLO CRYPTO";
    caesar_encrypt(msg, 3);
    printf("Encrypted: %s\n", msg);
    caesar_decrypt(msg, 3);
    printf("Decrypted: %s\n", msg);
    return 0;
}
\end{lstlisting}

\subsection{C++: Cipher Caesar}
\begin{lstlisting}[style=cppstyle,caption={Implementasi Caesar Cipher dalam C++}]
#include <iostream>
#include <string>
#include <cctype>

std::string caesar_encrypt(std::string plaintext, int shift) {
    std::string ciphertext = "";
    for (char c : plaintext) {
        if (isalpha(c)) {
            char base = isupper(c) ? 'A' : 'a';
            ciphertext += base + (c - base + shift + 26) % 26;
        } else {
            ciphertext += c;
        }
    }
    return ciphertext;
}

std::string caesar_decrypt(std::string ciphertext, int shift) {
    return caesar_encrypt(ciphertext, 26 - shift);
}

int main() {
    std::string msg = "HELLO CRYPTO";
    std::cout << "Encrypted: " << caesar_encrypt(msg, 3) << std::endl;
    std::cout << "Decrypted: " << caesar_decrypt(
        caesar_encrypt(msg, 3), 3) << std::endl;
    return 0;
}
\end{lstlisting}

\subsection{Pascal: Cipher Caesar}
\begin{lstlisting}[style=pascalstyle,caption={Implementasi Caesar Cipher dalam Pascal}]
program CaesarCipher;

var
  plaintext, ciphertext: string;
  shift, i: integer;
  c: char;

function caesar_encrypt(s: string; k: integer): string;
var
  result: string;
  i: integer;
  base: char;
begin
  result := '';
  for i := 1 to length(s) do
  begin
    if (s[i] >= 'A') and (s[i] <= 'Z') then
    begin
      base := 'A';
      result := result + chr(ord(base) + (ord(s[i]) - ord(base) + k) mod 26);
    end
    else if (s[i] >= 'a') and (s[i] <= 'z') then
    begin
      base := 'a';
      result := result + chr(ord(base) + (ord(s[i]) - ord(base) + k) mod 26);
    end
    else
      result := result + s[i];
  end;
  caesar_encrypt := result;
end;

function caesar_decrypt(s: string; k: integer): string;
begin
  caesar_decrypt := caesar_encrypt(s, 26 - k);
end;

begin
  plaintext := 'HELLO CRYPTO';
  shift := 3;
  ciphertext := caesar_encrypt(plaintext, shift);
  writeln('Encrypted: ', ciphertext);
  writeln('Decrypted: ', caesar_decrypt(ciphertext, shift));
end.
\end{lstlisting}

\subsection{MATLAB: Algoritma Euclid Diperluas}
\begin{lstlisting}[style=matlab,caption={Extended Euclidean Algorithm dalam MATLAB}]
function [g, x, y] = extended_gcd(a, b)
% Extended Euclidean Algorithm: ax + by = gcd(a,b)
    if b == 0
        g = a;
        x = 1;
        y = 0;
    else
        [g, x1, y1] = extended_gcd(b, mod(a, b));
        x = y1;
        y = x1 - floor(a/b) * y1;
    end
end

% Contoh: invers modular dari a mod m
function inv = mod_inverse(a, m)
    [g, x, ~] = extended_gcd(mod(a, m), m);
    if g ~= 1
        inv = [];  % tidak ada invers
    else
        inv = mod(x, m);
    end
end
\end{lstlisting}

\subsection{C: Modular Exponentiation (Fast Power)}
\begin{lstlisting}[style=cstyle,caption={Fast Modular Exponentiation dalam C}]
#include <stdio.h>

long long mod_pow(long long base, long long exp, long long mod) {
    long long result = 1;
    base = base % mod;
    while (exp > 0) {
        if (exp % 2 == 1)
            result = (result * base) % mod;
        exp = exp / 2;
        base = (base * base) % mod;
    }
    return result;
}

int main() {
    printf("2^10 mod 1000 = %lld\n", mod_pow(2, 10, 1000));
    printf("3^644 mod 645 (RSA) = %lld\n", mod_pow(3, 644, 645));
    return 0;
}
\end{lstlisting}

%==============================================================================
\section{Daftar Referensi}
%==============================================================================

\subsection{Buku Utama}
\begin{enumerate}
\item Stinson, D. R., \& Paterson, M. (2018). \textit{Cryptography: Theory and Practice} (4th ed.). CRC Press.
\item Katz, J., \& Lindell, Y. (2020). \textit{Introduction to Modern Cryptography} (3rd ed.). CRC Press.
\item Trappe, W., \& Washington, L. C. (2006). \textit{Introduction to Cryptography with Coding Theory} (2nd ed.). Pearson.
\item Schneier, B. (1996). \textit{Applied Cryptography} (2nd ed.). John Wiley \& Sons.
\end{enumerate}

\subsection{Buku Pendukung}
\begin{enumerate}
\item Stallings, W. (2017). \textit{Cryptography and Network Security: Principles and Practice} (7th ed.). Pearson.
\item Ferguson, N., Schneier, B., \& Kohno, T. (2010). \textit{Cryptography Engineering}. Wiley.
\item Menezes, A. J., van Oorschot, P. C., \& Vanstone, S. A. (1996). \textit{Handbook of Applied Cryptography}. CRC Press.
\end{enumerate}

\subsection{Sumber Online}
\begin{enumerate}
\item NIST Cryptographic Standards: \url{https://csrc.nist.gov/}
\item Cryptography Stack Exchange: \url{https://crypto.stackexchange.com/}
\end{enumerate}

%==============================================================================
\section{Pengesahan}
%==============================================================================

\begin{table}[!htbp]
\centering
\small
\begin{tabular}{|>{\centering\arraybackslash}p{0.3\linewidth}|>{\centering\arraybackslash}p{0.3\linewidth}|>{\centering\arraybackslash}p{0.3\linewidth}|}
\hline
\textbf{Disiapkan Oleh} & \textbf{Diperiksa Oleh} & \textbf{Disahkan Oleh} \\
\textbf{(Penyusun Silabus)} & \textbf{(Ketua Prodi)} & \textbf{(Dekan)} \\
\hline
Ketua Tim Penyusun Kurikulum / Penyusun Silabus & & \\
\hline
 & & \\[1.2cm]
Nama: \underline{\hspace{2.8cm}} & Nama: \underline{\hspace{2.8cm}} & Nama: \underline{\hspace{2.8cm}} \\
Tanggal: \underline{\hspace{2.8cm}} & Tanggal: \underline{\hspace{2.8cm}} & Tanggal: \underline{\hspace{2.8cm}} \\
\hline
\end{tabular}
\caption{Pengesahan Silabus Mata Kuliah Kriptografi (format BAN-PT/LAM)}
\end{table}

\end{document}
